\section{Look-ahead more than doubles the gain}
\label{sec:reward}

\begin{table*}[t]
\centering
\small
\begin{tabular}{lrrrrrr}
\toprule
& \multicolumn{3}{c}{\primaryjudge\ (training oracle)} & \multicolumn{3}{c}{\heldoutjudge\ (held out)}\\
\cmidrule(lr){2-4}\cmidrule(lr){5-7}
Instrument & $\Delta$ & \dz & $p_{\text{holm}}$ & $\Delta$ & \dz & $p_{\text{holm}}$\\
\midrule
Q1+Q2 (reward) & $+0.765$ & $0.905$ & $<.001$ & $+0.616$ & $1.030$ & $<.001$\\
\quad Q1        & $+0.858$ & $0.902$ & $<.001$ & $+0.865$ & $1.152$ & $<.001$\\
\quad Q2        & $+0.671$ & $0.864$ & $<.001$ & $+0.367$ & $0.609$ & $<.001$\\
WAI-SR          & $+0.291$ & $0.513$ & $<.001$ & $+0.288$ & $0.442$ & $<.001$\\
CSQ-8           & $+0.289$ & $0.482$ & $<.001$ & $+0.451$ & $0.678$ & $<.001$\\
MI-SAT          & $+0.352$ & $0.531$ & $<.001$ & $+0.503$ & $0.829$ & $<.001$\\
MITI            & $+0.615$ & $0.735$ & $<.001$ & $+0.276$ & $0.502$ & $<.001$\\
PCT             & $+0.111$ & $0.516$ & $<.001$ & $+0.113$ & $0.563$ & $<.001$\\
MICI ($\downarrow$) & $-0.627$ & $-1.862$ & $<.001$ & $-0.422$ & $-1.567$ & $<.001$\\
\bottomrule
\end{tabular}
\caption{\textbf{$\Kf$ minus $\Kz$ at the matched iteration-10 endpoint}, persona-paired over the
96 shared personas ($n{=}96$). A \emph{positive} $\Delta$ means look-ahead scored higher; MICI is
lower-is-better, so its negative $\Delta$ also favours $\Kf$. Look-ahead wins on every instrument
under both graders. Holm correction is across the nine rows within a grader.
Source: \texttt{k\_endpoints} (lookahead/reward family).}
\label{tab:endpoint}
\end{table*}

\subsection{The headline}

At the matched iteration-10 endpoint, look-ahead wins on \textbf{every instrument under both
graders} (Table~\ref{tab:endpoint}). On the rewarded rubric the gap is $+0.765$ ($\dz\,0.905$)
under the training oracle and $+0.616$ ($\dz\,1.030$) under the held-out judge. That the held-out
judge --- a different model family that never played the patient and never touched training ---
reports a \emph{larger} standardised effect than the grader that was optimised against is the
single most important fact in this section: the gain is not an artifact of the optimisation target.

Measured against each arm's own base, look-ahead more than doubles what the same optimizer extracts
from the same oracle: $+1.554$ versus $+0.686$ on Q1+Q2 under the training oracle
($1.554 / 0.686 = 2.27\times$), and $+1.038$ versus $+0.396$ under the held-out judge
($1.038 / 0.396 = 2.62\times$). In absolute terms the $\Kf$ policy ends at $4.517$ on a 1--5 scale
against the $\Kz$ policy's $3.753$, from statistically indistinguishable bases ($2.963$ and
$3.067$; the base-vs-base contrast is $+0.104$, $\dz\,0.115$, n.s.).

\subsection{The effect has an onset}

\begin{figure*}[t]
\centering
\includegraphics[width=\textwidth]{figures/k_headline_q1q2_grpo.png}
\caption{The two GRPO arms only. \textbf{Top:} Q1+Q2 by iteration under each grader, each arm's own
base dotted. \textbf{Bottom:} the persona-paired $\Kf - \Kz$ contrast, so \emph{up is look-ahead
winning} --- note this is the opposite sign to the underlying tables, which report $\Kz - \Kf$.
Iteration 0 is the two arms' independent base draws (the noise floor); the shaded band is the
$\pm 0.10$ oracle repeatability interval. The two graders' panels have independent $y$-axes and are
not on a common scale. Bars are persona-bootstrap 95\% CIs; filled markers cleared Holm.
The four-arm version of this figure, showing the PTO arms of the parent experiment for context, is
Figure~\ref{fig:fourarm-traj} in the appendix.
Source: \texttt{k\_headline\_q1q2\_grpo} (lookahead/reward family).}
\label{fig:headline}
\end{figure*}

Look-ahead is not free from the first step, and the trajectory matters for anyone deciding how long
to train. For the first three iterations the two arms are statistically indistinguishable on Q1+Q2
under both graders. The advantage opens at iteration 4 and is Holm-significant at 6 of the 10
iterations under each grader (primary: 4, 6, 7, 8, 9, 10; held-out: 4, 5, 6, 7, 9, 10), and it is
\emph{largest at the endpoint} --- there is no sign of it having peaked
(Figure~\ref{fig:headline}).

Two readings are ruled out by the same figure. This is not a case of $\Kz$ collapsing while $\Kf$
holds steady: $\Kf$ climbs throughout. Nor is it a single-iteration artifact --- the gap is
Holm-significant at every iteration from 6 to 10 under the training oracle. It is not smooth,
however, and we do not describe it as such: under the training oracle it runs
$+0.263, +0.196, +0.172, +0.647, +0.765$ over iterations 6--10, so most of the final margin opens
between iterations 8 and 9. Part of that late jump is the $\Kz$ arm falling rather than the $\Kf$
arm rising --- $\Kz$ scores $4.082$ at its own iteration-8 peak and $3.753$ at iteration 10 --- and
§\ref{sec:behaviour} identifies what the $\Kz$ policy was doing over exactly those iterations.
\todo{once the replicate draw of GRPO\_LA5@10 exists, state whether the endpoint survives a second
96-conversation sample; the iteration-9 dip in the $\Kz$ arm is a single draw}

\subsection{What the gain is made of}

The two components of the reward move differently, and the held-out judge localises the effect
more sharply than the trained-against grader does. Q1 (session satisfaction) carries the larger
share under both graders ($+0.858$ / $+0.865$), while Q2 (working alliance) is smaller and grader
dependent ($+0.671$ primary, $+0.367$ held out). The instruments \emph{outside} the reward move
too, which is what distinguishes a genuine gain from a target-specific one: MI-SAT
($\dz\,0.531$ / $0.829$), CSQ-8 ($0.482$ / $0.678$), WAI-SR ($0.513$ / $0.442$), the MITI-style
technique coder ($0.735$ / $0.502$), and patient change talk ($0.516$ / $0.563$) all favour
look-ahead under both graders. MI-inconsistent behaviour, the one lower-is-better axis, shows the
largest standardised effect in the table ($\dz\,-1.862$ primary, $-1.567$ held out); §\ref{sec:behaviour}
takes that apart, because it is where the two graders stop agreeing about magnitude.
